Language-conditioned robot manipulation is increasingly evaluated with large policy models, simulation suites, and embodied-agent benchmarks \cite{shridhar2022cliport,jiang2023vima,brohan2023rt1,brohan2023rt2,openx2023}. These benchmarks are essential for measuring end-to-end behavior, but aggregate success rates can obscure where failures enter the perception-to-action pipeline. A system may fail because it selects the wrong object from a language query, because a correct detection is lifted to an unstable 3D target, because depth is corrupted by occlusion or sparsity, or because a downstream executor cannot tolerate a small target displacement.

EmbodiedStressBench studies this intermediate language-to-target stage directly. The central unit of analysis is a diagnostic episode: a task, language query, detection list, selector, RGB-D target source, stressor family, stress level, and seed. The benchmark reports diagnostic target success, target error, failure type, and oracle gap, where the oracle target isolates task feasibility from non-oracle target generation. This scope is narrower than a full manipulation policy benchmark by design. It is intended to complement policy-level suites such as ManiSkill, CALVIN, SIMPLER, and VLABench \cite{gu2021maniskill,gu2023maniskill2,mees2022calvin,li2024simpler,zhang2024vlabench}.

The Scientific Reports version is framed as a reproducible simulation study. We ask five evidence-driven questions: whether controlled stressors expose target-generation failures, whether RGB-D target sources trade precision for robustness, whether detection-list semantic ambiguity creates measurable selection failures, whether learned open-vocabulary detectors can be evaluated through the same protocol, and whether diagnostic target success is currently calibrated to scripted execution success. The last question is answered negatively by the present executor smoke test unless the oracle execution gate is passed by a subsequent calibration run, and we treat negative execution evidence as a boundary on the claims rather than as a discarded experiment.

Our contributions are: (i) a parameterized stressor taxonomy for detection-list ambiguity, visual, geometric, and target-displacement perturbations; (ii) an oracle-gap and failure-taxonomy protocol for language-to-target generation; (iii) large-scale seed-controlled diagnostic evidence with confidence intervals and threshold sensitivity; (iv) a detector bridge that evaluates metadata-assisted, first-detection, and GroundingDINO selectors through the same RGB-D lifting interface; and (v) an explicit result-to-claim audit that distinguishes supported diagnostic findings from unresolved execution calibration.
